\documentclass[aspectratio=169]{beamer}

%\includeonlyframes{current}

\usepackage[utf8]{inputenc}
\usepackage[american]{babel}
\usepackage{amsmath,amsthm}
\usepackage{unicode}
\usepackage{array,tabularx}
\usepackage{ifthen}
\usepackage{tikz}
\usetikzlibrary{matrix,decorations,decorations.text,calc,arrows,snakes,shapes,positioning}
\usepackage{tikzsymbols}
\usepackage[backend=biber,citestyle=authoryear-comp,bibstyle=beamer,doi=false,isbn=false,url=true,maxnames=10]{biblatex}
\bibliography{refs}

\usepackage{ulem}

\mode<presentation>{%
  \usetheme{ibm}
}

\newcommand{\C}{ℂ}
\newcommand{\R}{ℝ}
\newcommand{\Z}{ℤ}
\newcommand{\N}{ℕ}
\newcommand{\Q}{ℚ}
\newcommand{\F}{\mathbb{F}}
\renewcommand{\P}{\mathbb{P}}
\renewcommand{\O}{\mathcal{O}}
\newcommand{\tildO}{\mathcal{\tilde{O}}}
\newcommand{\poly}{\operatorname{poly}}
\newcommand{\polylog}{\operatorname{polylog}}
\newcommand{\rank}{\operatorname{rank}}
\newcommand{\End}{\operatorname{End}}
\newcommand{\Hom}{\operatorname{Hom}}
\newcommand{\Gal}{\operatorname{Gal}}
\newcommand{\chr}{\operatorname{char}}
\newcommand{\Cl}{\operatorname{Cl}}
\newcommand{\GL}{\operatorname{GL}}
\renewcommand{\a}{{\mathfrak{a}}}
\renewcommand{\b}{{\mathfrak{b}}}
\newcommand{\p}{{\mathfrak{p}}}
\newcommand{\q}{{\mathfrak{q}}}
\newcommand{\g}{{\mathfrak{g}}}
\newcommand{\G}{{\mathcal{G}}}
\newcommand{\E}{{\mathcal{E}}}
\newcommand{\cyc}[1]{{〈 #1 〉}}
\newcommand{\ord}{\operatorname{ord}}
\newcommand{\mat}[1]{\left(\begin{smallmatrix}#1\end{smallmatrix}\right)}
\newcommand{\from}{\overset{\$}{\leftarrow}}

\newcommand{\bl}[1]{\textcolor{blue}{#1}}
\newcommand{\rd}[1]{\textcolor{red}{#1}}
\newcommand{\gr}[1]{\textcolor{green}{#1}}
\newcommand{\og}[1]{\textcolor{orange}{#1}}

\definecolor{light blue}{RGB}{0,102,204}

\newcommand{\myedge}[3]{
  \draw[#3] (360/\crater*#1 : \diam) to[bend right] (360/\crater*#2 : \diam);
}
\newcommand{\sk}[4]{
  \draw[very thick,blue]   (0,0) -- (0,#1);
  \draw[very thick,red]    (1,0) -- (1,#2);
  \draw[very thick,green]  (2,0) -- (2,#3);
  \draw[very thick,orange] (3,0) -- (3,#4);
}
\newcommand{\axes}[4]{
  \clip (#1,#3) rectangle (#2,#4);
  \draw [thin, gray, -latex] (#1,0) -- (#2,0);% Draw x axis
  \draw [thin, gray, -latex] (0,#3) -- (0,#4);% Draw y axis
}
\pgfkeys{/lattice/.code n args={4}{\tikzset{cm={#1,#2,#3,#4,(0,0)}}}}
\newcommand{\lattice}[2]{
  \draw[style=help lines,dashed] (#1-1,#1-1) grid[step=1] (#2+1,#2+1);
  \foreach \x in {#1,...,#2}{
    \foreach \y in {#1,...,#2}{
      \node[draw,circle,inner sep=2pt,fill] at (\x,\y) {};
      % Places a dot at those points
    }
  }
}

\newenvironment<>{goodblock}[1]{%
  \begin{actionenv}#2%
      \def\insertblocktitle{#1}%
      \par%
      \mode<presentation>{%
        \setbeamercolor{block title}{fg=green!60!black,bg=green!50!white}
       \setbeamercolor{block body}{bg=green!20!white}
     }%
      \usebeamertemplate{block begin}}
    {\par\usebeamertemplate{block end}\end{actionenv}}
\newenvironment<>{mehblock}[1]{%
  \begin{actionenv}#2%
      \def\insertblocktitle{#1}%
      \par%
      \mode<presentation>{%
        \setbeamercolor{block title}{fg=yellow!50!black,bg=yellow!50!white}
       \setbeamercolor{block body}{bg=yellow!20!white}
     }%
      \usebeamertemplate{block begin}}
    {\par\usebeamertemplate{block end}\end{actionenv}}
\newenvironment<>{badblock}[1]{%
  \begin{actionenv}#2%
      \def\insertblocktitle{#1}%
      \par%
      \mode<presentation>{%
        \setbeamercolor{block title}{fg=red!60!black,bg=red!50!white}
       \setbeamercolor{block body}{bg=red!20!white}
     }%
      \usebeamertemplate{block begin}}
    {\par\usebeamertemplate{block end}\end{actionenv}}


../share/isographs.tex

\title{SQIsign}
\subtitle{The new herald of isogeny based crypto}
\author{Luca De Feo}
\date[December 15, 2023, CIFRIS23]{December 15, 2023, CIFRIS23}
\institute{IBM Research Zürich}

\begin{document}

\frame[plain]{\titlepage}

%%

\begin{frame}{Isogenies and signing}
  \begin{description}
  \item[2012]<1-> \textcolor{blue}{SIDH PoKs (DF, Jao, Plût), signatures (Yoo, Azarderakhsh, Jao)}
  \item[2017]<3-> \textcolor{red}{Galbraith--Petit--Silva}
  \item[2019]<2-> \textcolor{purple}{SeaSign (DF, Galbraith)}
  \item[]<2-> \textcolor{purple}{CSI-FiSh (Beullens, Kleinjung, Vercauteren)}
  \item[2020]<3-> \textcolor{red}{SQIsign (DF, Kohel, Leroux, Petit, Wesolowski)}
  \item[2022]<1-> \textcolor{blue}{Ternary SIDH PoKs (DF, Dobson, Galbraith, Zobernig)}
  \item[2023]<1-> \textcolor{blue}{\textit{Curves you can trust} \tiny(Basso, Codogni, Connolly, DF, Fouotsa, Lido, Morrison, Panny, Patranabis, Wesolowski)}
  \item[2024]<3-> \textcolor{red}{SQIsignHD (Dartois, Leroux, Robert, Wesolowski)}
  \end{description}

  \bigskip
  
  \newcolumntype{C}{>{\centering\arraybackslash}X}
  \begin{tabularx}{\textwidth}{ C | C | C }
    \color{blue} SIDH & \uncover<2->{\color{purple} CSIDH} & \uncover<3->{\color{red} Deuring} \\
    \uncover<1->{Best security} & \uncover<2->{Group, ring, threshold, \dots} & \uncover<3->{Most compact}
  \end{tabularx}
\end{frame}

%%

\begin{frame}{Elliptic curves}
  \begin{columns}
    \begin{column}{0.4\textwidth}
      {\Large
      \[y^2 = x^3 + ax + b\]
      }
      \begin{block}{Bezout's theorem}
        Every line cuts $E$ in exactly three points (counted with
        multiplicity).
      \end{block}

      Define a \emph{group law} such that any three colinear points
      add up to zero.
    \end{column}
    \begin{column}{0.6\textwidth}
      \begin{center}
        \begin{tikzpicture}[domain=-2.4566:4,samples=100,yscale=1/2]
          \draw plot (\x,{sqrt(\x*\x*\x-4*\x+5)});
          \draw plot (\x,{-sqrt(\x*\x*\x-4*\x+5)});

          \draw[thin,gray,-latex] (0,-7) -- (0,7);
          \draw[thin,gray,-latex] (-3,0) -- (4,0);
          \draw (-3,1) -- (4,8/3+3);
          \begin{scope}[every node/.style={draw,circle,inner sep=1pt,fill},cm={1,2/3,0,0,(0,3)}]
            \node at (-2.287980,0) {};
            \node at (-0.535051,0) {};
            \node at (3.267475,0) {};
          \end{scope}
          \begin{scope}[every node/.style={yshift=0.3cm},cm={1,2/3,0,0,(0,3)}]
            \node at (-2.287980,0) {$P$};
            \node at (-0.535051,0) {$Q$};
            \node at (3.267475,0) {$R$};
          \end{scope}

          \draw[dashed] (3.267475,3.267475*2/3+3) -- (3.267475,-3.267475*2/3-3) 
          node[draw,circle,inner sep=1pt,fill] {}
          node[xshift=-0.1cm,anchor=east] {$P+Q$};
        \end{tikzpicture}
      \end{center}
    \end{column}
  \end{columns}
\end{frame}

%%

\begin{frame}{$N$-torsion}
  \Large\centering
  
  \[E[N] \quad≃\quad ℤ/Nℤ × ℤ/Nℤ\]

  \bigskip
  \normalsize
  *over an algebraically closed field
\end{frame}

%% 

\begin{frame}
  \Large
  \begin{description}
    \setlength{\itemsep}{4em}
  \item[Isogenies =] finite-kernel group morphisms: \emph{$E \to E/K$}
  \item[Endomorphisms =] isogenies \emph{$E \to E$}
  \end{description}
\end{frame}

%%

\begin{frame}{Isogenies: an example over $\F_{11}$}
  \centering
  \begin{tikzpicture}[scale=0.4]
    \begin{scope}
      \node[anchor=center] at (0,7) {$E \;:\; y^2 = x^3 + x$};

      \uncover<-1>{
        \draw[thin,gray] (0,-6) -- (0,6);
        \draw[thin,gray] (-6,0) -- (6,0);
      }

      \foreach \x/\y in {0/0,5/3,-4/3,-3/5,-2/1,-1/3} {
        \draw[blue,fill] (\x,\y) circle (0.2) node(E_\x_\y){}
        (\x,-\y) circle (0.2) node(E_\x_-\y){};
      }

      \uncover<2->{\draw[red,fill] (0,0) circle (0.3);}
    \end{scope}

    \draw[black!10!white,thick] (10,-7) -- +(0,14);
    
    \begin{scope}[shift={(20,0)}]
      \node at (0,7) {$E' \;:\; y^2 = x^3 - 4x$};

      \uncover<-1>{
        \draw[thin,gray] (0,-6) -- (0,6);
        \draw[thin,gray] (-6,0) -- (6,0);
      }

      \foreach \x/\y in {0/0,2/0,3/2,4/2,6/4,-2/0,-1/5} {
        \draw[color=blue,fill] (\x,\y) circle (0.2) node(F_\x_\y){}
        (\x,-\y) circle (0.2) node(F_\x_-\y){};
      }
    \end{scope}

    \begin{scope}[color=red,-latex,dashed]
      \begin{uncoverenv}<2->
        \path
        (E_5_3) edge (F_3_2)
        (E_-4_3) edge (F_4_-2)
        (E_-3_5) edge (F_4_2)
        (E_-2_1) edge (F_3_-2)
        (E_-1_3) edge (F_-2_0);
      \end{uncoverenv}
      \begin{uncoverenv}<2->
        \path
        (E_5_-3) edge (F_3_-2)
        (E_-4_-3) edge (F_4_2)
        (E_-3_-5) edge (F_4_-2)
        (E_-2_-1) edge (F_3_2)
        (E_-1_-3) edge (F_-2_0);
      \end{uncoverenv}
    \end{scope}
  \end{tikzpicture}
  
  \begin{columns}
    \begin{column}{0.5\textwidth}
      \[\phi(x,y) = \left(\frac{x^2 + 1}{x},\quad y\frac{x^2-1}{x^2}\right)\]
    \end{column}
    \begin{column}{0.5\textwidth}
      \begin{itemize}
      \item<2-> Kernel generator in \alert{red}.
      \item<2-> This is a degree $2$ map.
      \item<2-> Analogous to $x\mapsto x^2$ in $\F_q^*$.
      \end{itemize}
    \end{column}
  \end{columns}
\end{frame}

%%

\begin{frame}{Some endomorphisms}
  \Large
  \begin{description}
    \setlength{\itemsep}{3em}
  \item[Scalar multiplication:] $[n] : E \to E/E[n]$
  \item[Frobenius:] $(x,y) \mapsto (x^p, y^p)$ \hfill(on curves over $\F_p$)
  \item[Automorphisms:] $(x,y) \mapsto (-x, iy)$ \hfill(on curve $y^2 = x^3 + x$)
  \end{description}
\end{frame}

%%

\begin{frame}{Endomorphisms = $2 \times 2$ matrices}
  Fix any basis \emph{$〈P,Q〉$} of \emph{$E[N]$}

  \begin{align*}
    \omega : E[N] &→ E[N]\\
    \begin{pmatrix}
      x\\y
    \end{pmatrix}
                  &↦
                    \begin{pmatrix}
                      a&b\\c&d
                    \end{pmatrix}
                              \begin{pmatrix}
                                x\\y
                              \end{pmatrix}
    \mod N
  \end{align*}
  
  \begin{block}{Tate's isogeny theorem}
    1-to-1 correspondence between endomorphisms and Galois-invariant
    maps \emph{$E[ℓ^∞] \to E[ℓ^∞]$},
    for any prime $ℓ ≠ p$.

    \[\emph{\End(E)⊗ℤ_ℓ \hookrightarrow \mathcal{M}_{2\times 2}(ℤ_ℓ)}\]
  \end{block}
\end{frame}

%% 

\begin{frame}{Endomorphisms = imaginary quadratic integers}
  \large
  Endomorphisms form a ring:
  \[ω ∘ (φ + ψ) \quad=\quad ω∘φ + ω∘ψ\]

  \bigskip
  Every endomorphism satisfies a quadratic equation
  \[ω^2 - tω + n = 0\]
  with $t,n ∈ ℤ$ and $t^2 - 4n ≤ 0$.
\end{frame}

%% 

\begin{frame}{Endomorphism rings}
  \Large $\End(E)$ is a free $ℤ$-module of rank \textcolor{blue}{1},
  \textcolor{purple}{2} or \textcolor{red}{4}. As a ring:

  \bigskip
  \begin{enumerate}
    \setlength{\itemsep}{1em}
  \item[\color{blue}1)] $\End(E) ≃ ℤ$;
  \item[\color{purple}2)] $\End(E)$ is isomorphic to an
    order\footnote[1]{order = subring of maximal rank} of a quadratic
    imaginary field;
  \item[\color{red}4)] $\End(E)$ is isomorphic to a maximal
    order\footnotemark[1] of the quaternion algebra ramified at $p$
    and $∞$.
  \end{enumerate}
\end{frame}

%% 

\begin{frame}{\color{purple}Quadratic imaginary number fields}
  \large
  \[ℚ(\sqrt{-D}) \qquad\text{with}\qquad D>0\]

  \bigskip
  \begin{itemize}
    \setlength{\itemsep}{1em}
  \item \textcolor{purple}{$2$-dimensional} $ℚ$-vector space,
  \item Unique maximal order $\quad=\quad$ ring of integers.
  \end{itemize}
\end{frame}

%%

\begin{frame}{\color{red}Quaternion algebras}
  (Assume $p=-1 \bmod 4$) The quaternion algebra \emph{$B_{p,∞}$} is:
  \begin{itemize}
  \item A \textcolor{red}{$4$-dimensional} $ℚ$-vector space with basis
    \emph{$(1,i,j,k)$}.
  \item A non-commutative \emph{division algebra}%
    \footnote{All elements have inverses.} %
    $B_{p,∞} = ℚ〈i,j〉$ with the relations:
    \[i^2 = -1, \quad j^2 = -p, \quad ij = -ji = k.\]
  \end{itemize}

  \begin{block}{Properties}
    \begin{itemize}
    \item All elements of $B_{p,∞}$ are \emph{quadratic algebraic
        numbers}.
    \item $B_{p,∞}⊗ℚ_ℓ≃\mathcal{M}_{2×2}(ℚ_ℓ)$ for all $ℓ≠p$.\\
    \item $B_{p,∞}⊗ℝ$ is isomorphic to Hamilton's quaternions.
    \item $B_{p,∞}⊗ℚ_p$ is a division algebra.
    \end{itemize}
  \end{block}
\end{frame}

%%

\begin{frame}{An example}
  The curve of $j$-invariant \emph{$1728$}
  \[E: y^2 = x^3 + x\]
  is supersingular over $\F_p$ iff $p=-1\mod 4$.

  \begin{block}{Endomorphisms}
    \emph{$\End(E) ⊂ ℚ〈ι,π〉$}, with:
    \begin{itemize}
    \item $π$ the Frobenius endomorphism, s.t. \emph{$π^2=-p$};
    \item $ι$ the map
      \[ι(x,y) = (-x,iy),\]
      where \emph{$i∈\F_{p^2}$} is a 4-th root of unity.
      Clearly, \emph{$ι^2=-1$}.
    \end{itemize}
    And \emph{$ιπ=-πι$}.
  \end{block}
\end{frame}

%%

\begin{frame}{Isogenies = finite subgroups}
  \large
  Let $K ⊂ E$ finite, there exists a unique (separable) isogeny
  \[\emph{\phi : E → E/K}\]
  up to post-composition by an isomorphism.

  \vspace{2em}

  $O(N)$ cyclic subgroups s.t. $\#K = N \quad\to\quad O(N)$ distinct isogenies
\end{frame}

%%

\begin{frame}{Isogenies = $2 × 2$ matrices}
  \large
  Denote by \emph{$\Hom(E, E')$} the group of isogenies $E → E'$

  \bigskip
  
  \begin{block}{Tate's isogeny theorem}
    \[\Hom(E, E') ⊗ ℤ_ℓ \quad≃\quad \Hom_k(E[ℓ^∞], E'[ℓ^∞]) \quad\hookrightarrow\quad \mathcal{M}_{2×2}(ℤ_ℓ)\]
  \end{block}

  \bigskip
  
  In particular, \emph{$\rank\Hom(E, E') = \rank\End(E)$}.
\end{frame}

%%

\begin{frame}{Isogenies = Ideals}
  \large\centering
  \begin{tikzpicture}
    \node (E1) at (0,0) {$E'$};
    \node (E) at (6,0) {$E$};

    \draw[-latex] (E) edge node[above] {$φ$} (E1);

    \uncover<2>{
      \draw[-latex] (E) edge[loop,out=45,in=-45,looseness=10] node[right] {$ω$} (E);
      \node at (3,-3) {$φ∘ω ∈ \Hom(E,E')$};
    }
    
    \uncover<3>{
      \draw[-latex] (E1) edge[loop,out=135,in=225,looseness=10] node[left] {$ω'$} (E1);
      \node at (3,-3) {$ω'∘φ ∈ \Hom(E,E')$};
    }
  \end{tikzpicture}
\end{frame}

%%

\begin{frame}{The Deuring correspondences}
  \centering\large
  \setlength{\tabcolsep}{2em}
  \renewcommand{\arraystretch}{1.8}
  \begin{tabular}{r l}
    \emph{Ellitpic curves} & \emph{Number fields / Quaternion algebras}\\
    \hline
    Endomorphisms & Algebraic integers\\
    Endomorphism ring & (Maximal) order\\
    Isogeny & Invertible ideal\\
    Isogeny degree & Ideal norm\\
    Isogeny composition & Ideal multiplication\\
    Isogenies \raisebox{-0.8em}{\tikz{\node (E) at (0,0) {$\bullet$}; \node (E1) at (2,0) {$\bullet$}; \draw[->] (E) edge[bend left] (E1) edge[bend right] (E1);}} & Ideal classes\\
    Dual isogeny & Conjugate ideal\\
  \end{tabular}
\end{frame}

%%
\begin{frame}{Two computational worlds}
  \centering
  \setlength{\tabcolsep}{2em}
  \renewcommand{\arraystretch}{1.5}
  \begin{tabular}{p{0.3\textwidth} c c}
    & \emph{Ordinary / CSIDH} & \emph{Supersingular}\\
    \hline
    $\rank\Hom(E,E')$ & 2 & 4\\
    Endomorphism algebra & number field & quaternion algebra\\
    Maximal orders & one & many \\
    Ideal class\dots & \dots group & \dots set\\
    Find isogeny $E → E'$ & \alert{hard} & \alert{hard}\\
    Convert isogenies $\leftrightarrow$ ideals & easy-ish & easy\\
    Compute $\End(E)$ & easy & \alert{hard}\\
  \end{tabular}  
\end{frame}

%%

\begin{frame}{The Deuring correspondence (for supersingular curves)}
  \begin{block}{An equivalence of categories (roughly)}
    \centering
    \begin{tikzpicture}
      \node (O) at (0,0) {$\O$};
      \node (O1) at (6,0) {$\O'$};
      \node (E) at (0,-1) {$E$};
      \node (E1) at (6,-1) {$E'$};
      
      \begin{scope}[gray,anchor=north]
        \node (Oc) at (-2,1) {left order};
        \node at (-2, 1.5) {$\{ω\in B_{p,\infty}\;|\; ωI=I\}$};
        \node (O1c) at (8,1) {right order};
        \node at (8,1.5) {$\{ω\in B_{p,\infty}\;|\; Iω=I\}$};
        \node (ac) at (3,1.5) {connecting ideal};
        
        \node (Ec) at (-2,-2) {supersingular curve};
        \node (E1c) at (8,-2) {supersingular curve};
        \node (phic) at (3,-2) {isogeny};
      \end{scope}
      
      \draw[->] (O) edge node[auto] (a) {$I$} (O1)
      (E) edge node[auto,swap] (phi) {$\phi_I$} (E1);
      \draw[dashed,->] (Oc) edge (O) (O1c) edge (O1) (ac) edge (a)
      (Ec) edge (E) (E1c) edge (E1) (phic) edge (phi);
    \end{tikzpicture}
  \end{block}
\end{frame}

%%

\begin{frame}{The effective Deuring correspondence}
  \centering\large
  \setlength{\tabcolsep}{2em}
  \renewcommand{\arraystretch}{1.8}
  \begin{tabular}{c c c}
    Input & Output \\
    \hline
    random $E$ & $\End(E)$ & \alert{hard}\\
    random $E$ & $ω ∈ \End(E)$ & \alert{hard}\\
    random $E$ & $\phi : E_0\to E$ & \alert{hard}\\
    \hline
    $\End(E)$ & $E$ & easy\\
    $\End(E)$, $\End(E')$ & connecting ideal & easy\\
    \hline
    $I ⊂ \End(E)$ & $\phi_I : E → E'$ & easy\\
    $\End(E)$,\;\; $\phi: E → E'$ & $I_\phi ⊂ \End(E)$,\;\; $\End(E')$ & easy
  \end{tabular}
\end{frame}

%%

\begin{frame}{A loose analogy: Signing based on factoring}\
  {\Large \[N = pq\]}

  \bigskip
  \begin{center}
    \begin{tabular}{l | c | c}
      & $\Z/N\Z$ & $\Z/p\Z \times \Z/q\Z$\\
      \hline
      multiplication & easy & easy\\
      inversion & easy & easy\\
      square roots & \alert{hard} & easy\\
      $n$-th roots & \alert{hard} & easy\\
    \end{tabular}
  \end{center}

  \begin{block}{Rabin's signature}
    \begin{description}
    \item[Sign:] $s \gets \sqrt{H(m; r)} \mod N$,
    \item[Verify:] $s^2 \overset{?}{=} H(m; r) \mod N$.
    \end{description}
  \end{block}
\end{frame}

%%

\begin{frame}{The endomorphism ring problem}
  \Large\centering
  \vfill
  Given a random supersingular curve $E$, compute $\End(E)$
  \vfill
\end{frame}

%%

\begin{frame}{Contagious knowledge}
  \centering
  \begin{tikzpicture}
    \pgfmathsetseed{12345}
    \foreach \i in {1,...,100} {
      \pgfmathparse{16*random()}
      \let\x\pgfmathresult
      \pgfmathparse{7*random()}
      \let\y\pgfmathresult
      \fill[black!20!white] (\x,\y) circle (1pt);
    }
    \foreach \i in {1,...,10} {
      \pgfmathparse{floor(15*random())}
      \let\x\pgfmathresult
      \pgfmathparse{floor(6*random())}
      \let\y\pgfmathresult
      \fill (\x,\y) circle (2pt);
      \uncover<2->{
        \foreach \rho in {0,1,2} {
          \draw[fill,-latex] (\x,\y) -- +(120*\rho:.3) circle (2pt);
          \uncover<3->{
            \foreach \sigma in {-1,1} {
              \draw[fill,-latex] (\x,\y) ++(120*\rho:.3) -- ++(120*\rho+60*\sigma:.3) circle (2pt);
            }
          }
        }
      }
    }
  \end{tikzpicture}
\end{frame}

%%

\begin{frame}{Galbraith--Petit--Silva (2017)}
  \large
  $\Sigma$-protocol with binary challenge:

  \bigskip
  \centering
  \begin{tikzpicture}
    \node (E0) at (1,3) {$E_0$};
    \node (EA) at (0,0) {$E_A$};
    \draw [blue,dashed,->] (E0) edge (EA);
    
    \uncover<2->{
    \node (Ec) at (5,0) {$E_c$};
      \draw [blue] [->] (EA) to (Ec);
    }
    \uncover<3->{
      \draw [blue,very thick] [->] (E0) -- (Ec);
    }
    
    \matrix [right] at (6,2) {
      \draw[dashed] (0,0) edge[blue] (1,0) (1,0) node {secret key};\\
      \uncover<2->{\draw (0,0) edge[blue] (1,0) (1,0) node {commitment};}\\
      \uncover<3->{\draw (0,0) edge[blue] (1,0)
        (0,0.2) edge[blue,very thick] (1,0.2)
        (1,0.1) node {response};}\\
    };
  \end{tikzpicture}  
\end{frame}

%%

\begin{frame}{Equivalent isogenies $\leftrightarrow$ equivalent ideals (KLPT)}
  \large\centering
  \begin{tikzpicture}
    \node (E) at (0,0) {$E$};
    \node (E1) at (6,0) {$E'$};
    \uncover<3->{
      \node (O) at (7,-3) {$\End(E)$};
      \node (O1) at (13,-3) {$\End(E')$};
    }

    \begin{scope}[-latex]
      \draw (E) edge[bend left] node[above] {$φ$} (E1);
      \uncover<2>{\draw (E) edge[bend right,dashed,red] node[red] {$×$} (E1);}
      \uncover<5->{\draw (E) edge[bend right] node[below] {$φ_J$} (E1);}
      \uncover<3->{\draw (O) edge[bend left] node[above] {$I_φ$} (O1);}
      \uncover<4->{\draw (O) edge[bend right] node[below] {$J = I_φα$} (O1);}
    \end{scope}
  \end{tikzpicture}
\end{frame}

%% 

\begin{frame}{SQIsign: Signatures from the effective Deuring correspondence}
  \begin{center}
    \begin{tikzpicture}
      \node (E0) at (1,2.5) {$E_0$};
      \uncover<2->{
        \node (E1) at (4,2.5) {$E_1$};
        \node (B) at (2.5,2.75) {$\psi$};
        \draw [blue] [->] (E0) to (E1);
      }
      \uncover<3->{
        \node (E2) at (4,1) {$E_2$};
        \node (A) at (4.25,1.75) {$\varphi$};
        \draw [red] [->] (E1) to (E2);
      }
      \node (EA) at (0,1) {$E_A$};
      \node (A) at (0.25,1.75) {$\tau$};
      \draw [blue,dashed] [->] (E0) to (EA);

      \uncover<4->{
        \node (B) at (2,1.25) {$\sigma$};
        \draw [blue,very thick] [->] (EA) -- (E2);
      }

      \matrix [right] at (6,2) {
        \node[] (l1) {}; \node (l2) [right of = l1, node distance=0.5in,label=right: commitment isogeny (prover)] {}; \draw [blue] [->] (l1) -- (l2); \\
        \node[] (l3) {}; \node (l4) [right of = l3, node distance=0.5in,label=right:challenge isogeny (verifier)] {}; \draw [red] [->] (l3) -- (l4); \\
        \node[] (l1) {}; \node (l2) [right of = l1, node distance=0.5in,label=right: response isogeny (prover)] {}; \draw [blue,very thick] [->] (l1) -- (l2); \\
        \node[] (l1) {}; \node (l2) [right of = l1, node distance=0.5in,label=right: secret key isogeny] {}; \draw  [dashed,blue] [-] (l1) -- (l2); \\
      };
    \end{tikzpicture}
  \end{center}

  \bigskip
  
  \emph{Most compact PQ signature scheme}: PK + Signature combined
  \textbf{5$\times$smaller} than Falcon.

  \begin{table}[h]
    \centering
    \begin{tabular}{ c c c c}
      Secret Key (bytes) & Public Key (bytes) & Signature (bytes) & Security \\
      \hline
      782 & 64 & 177 & NIST-1 \\
      1138 & 96 & 263 & NIST-3 \\
      1509 & 128 & 335 & NIST-5 \\      
      \hline
    \end{tabular}
  \end{table}
\end{frame}

%% 

\begin{frame}[plain]
  \centering
  \begin{tikzpicture}[remember picture,overlay]
    \begin{scope}[xscale=1.7,yshift=-15,opacity=0.8]
      \def\crater{12}
      \def\jumpa{-8}
      \def\jumpb{9}
      \def\diam{5cm}

      \foreach \i in {1,...,\crater} {
        \draw[blue] (360/\crater*\i : \diam) to[bend right] (360/\crater*\i+360/\crater : \diam);
        \draw[red] (360/\crater*\i : \diam) to[bend right] (360/\crater*\i+\jumpa*360/\crater : \diam);
        \draw[green] (360/\crater*\i : \diam) to[bend right=50] (360/\crater*\i+\jumpb*360/\crater : \diam);
      }
    \end{scope}
    
    
    \draw (0,0.5) node{\Huge\bf Thank you};
    \draw (0,-0.6) node{\large\url{https://defeo.lu/}};
    \draw (0,-1.3) node{\large\includegraphics[height=0.9em]{mastodon.png}~\href{https://twitter.com/luca_defeo}{@luca\_defeo@ioc.exchange}};
    \draw (0,-1.9) node{\large\includegraphics[height=0.9em]{twitter.png}~\href{https://twitter.com/luca_defeo}{@luca\_defeo}};
  \end{tikzpicture}
\end{frame}

%%

\end{document}


% LocalWords:  Isogeny abelian isogenies hyperelliptic supersingular Frobenius
% LocalWords:  isogenous
